% SPDX-License-Identifier: MIT
% Copyright (c) 2026 Evgeny Baulin
% Optimization Methods, Seminar 09: theory handout (Russian).
% Generated by checks/09. Linear Programming as the Foundation of Integer Optimization.ipynb from web/09/js/content-*.js and content/students.json. Do not edit by hand.
% Compile with: latexmk -xelatex -interaction=nonstopmode -halt-on-error Theory_ru.tex
\documentclass[11pt,a4paper]{article}

\usepackage{fontspec}
\setmainfont{cmunrm.otf}[BoldFont=cmunbx.otf, ItalicFont=cmunti.otf, BoldItalicFont=cmunbi.otf]
\setsansfont{cmunss.otf}[BoldFont=cmunsx.otf, ItalicFont=cmunsi.otf, BoldItalicFont=cmunso.otf]
\setmonofont{cmuntt.otf}[BoldFont=cmuntb.otf, ItalicFont=cmunit.otf, BoldItalicFont=cmuntx.otf]
\usepackage{polyglossia}
\setdefaultlanguage{russian}
\setotherlanguage{english}
\newfontfamily\cyrillicfont{cmunrm.otf}[Ligatures=TeX, BoldFont=cmunbx.otf, ItalicFont=cmunti.otf, BoldItalicFont=cmunbi.otf]
\newfontfamily\cyrillicfontsf{cmunss.otf}[Ligatures=TeX, BoldFont=cmunsx.otf, ItalicFont=cmunsi.otf, BoldItalicFont=cmunso.otf]
\newfontfamily\cyrillicfonttt{cmuntt.otf}[BoldFont=cmuntb.otf, ItalicFont=cmunit.otf, BoldItalicFont=cmuntx.otf]

\usepackage[a4paper,left=2.3cm,right=2.3cm,top=2.4cm,bottom=2.4cm,headheight=14pt]{geometry}
\usepackage{amsmath,amssymb}
\usepackage{graphicx}
\usepackage{array,booktabs,longtable}
\usepackage{enumitem}
\usepackage{xcolor}
\usepackage[most]{tcolorbox}
\usepackage{titlesec}
\usepackage{fancyhdr}
\usepackage{needspace}
\usepackage[expansion=false]{microtype}
\usepackage[colorlinks=true,linkcolor=hseblue,urlcolor=hseblue,citecolor=hseblue]{hyperref}

\definecolor{hseblue}{RGB}{16,45,105}
\definecolor{hseaccent}{RGB}{198,40,40}
\definecolor{hsegray}{RGB}{108,114,126}
\definecolor{hselight}{RGB}{240,243,248}
\definecolor{hsewarm}{RGB}{253,244,232}

\hypersetup{pdftitle={Методы оптимизации. Семинар 9. Конспект теории}, pdfauthor={Евгений Баулин}}

\titleformat{\section}{\color{hseblue}\Large\bfseries}{\thesection}{0.6em}{}
\titleformat{\subsection}{\color{hseblue}\large\bfseries}{\thesubsection}{0.6em}{}
\titlespacing*{\section}{0pt}{16pt plus 4pt minus 2pt}{8pt plus 2pt}
\titlespacing*{\subsection}{0pt}{12pt plus 3pt minus 2pt}{6pt plus 2pt}
\pretocmd{\section}{\Needspace{7\baselineskip}}{}{\errmessage{cannot patch section}}
\pretocmd{\subsection}{\Needspace{5\baselineskip}}{}{\errmessage{cannot patch subsection}}

\tcbset{handoutbox/.style={enhanced, breakable, sharp corners, boxrule=0pt, left=7pt, right=7pt, top=5pt, bottom=5pt, before skip=8pt plus 2pt, after skip=8pt plus 2pt}}
\newtcolorbox{examplebox}{handoutbox, colback=hselight, colframe=hselight, before upper={\textcolor{hseblue}{\bfseries Пример.}\ }}
\newtcolorbox{whybox}{handoutbox, colback=white, colframe=hseblue, leftrule=2.5pt, top=3pt, bottom=3pt, before upper={\textcolor{hseaccent}{\bfseries Почему.}\ }}
\newtcolorbox{answerbox}{handoutbox, colback=hsewarm, colframe=hsewarm, before upper={\textcolor{hseaccent}{\bfseries Ответ.}\ }}
\newtcolorbox{conventionbox}{handoutbox, colback=hselight, colframe=hseblue, leftrule=2.5pt, before upper={\textcolor{hseblue}{\bfseries Обозначения.}\par\smallskip}}

\newcommand{\handoutfigure}[3]{\par\medskip\noindent\begin{minipage}{\linewidth}\centering
  \includegraphics[width=#1\linewidth]{figures/#2.pdf}\par\smallskip
  \begin{minipage}{0.94\linewidth}\small\textcolor{hseblue}{\bfseries Рисунок.}\ #3\end{minipage}\end{minipage}\par\medskip}
\newcommand{\problemhead}[3]{\par\bigskip\noindent{\color{hseblue}\large\bfseries #1.\ #2}\par\nopagebreak
  \noindent{\small\itshape\color{hsegray}#3}\par\nopagebreak\smallskip}
\newtcolorbox{definitionbox}[1]{handoutbox, colback=white, colframe=hseblue, leftrule=2.5pt, before upper={\textcolor{hseblue}{\bfseries #1}\ }}
\newtcolorbox{theorembox}[1]{handoutbox, colback=hselight, colframe=hseblue, leftrule=2.5pt, before upper={\textcolor{hseblue}{\bfseries #1}\ }}
\newtcolorbox{proofbox}{handoutbox, colback=white, colframe=hsegray, leftrule=1.2pt, top=3pt, bottom=3pt, before upper={\textcolor{hsegray}{\bfseries\itshape Доказательство.}\ }}
\newtcolorbox{notebox}[1]{handoutbox, colback=hselight, colframe=hselight, before upper={\textcolor{hseblue}{\bfseries #1}\par\smallskip}}
\newtcolorbox{trapbox}{handoutbox, colback=hsewarm, colframe=hseaccent, leftrule=2.5pt, before upper={\textcolor{hseaccent}{\bfseries Ловушка.}\ }}
\newtcolorbox{selfcheckbox}{handoutbox, colback=white, colframe=hseaccent, leftrule=2.5pt, top=3pt, bottom=3pt, before upper={\textcolor{hseaccent}{\bfseries Проверьте себя.}\ }}
\newcommand{\handoutdisplay}[1]{\sbox0{$\displaystyle #1$}\ifdim\wd0>\linewidth
  \par{\centering\binoppenalty=10000\relpenalty=0\relax$\displaystyle #1$\par}\else\[#1\]\fi}
\newcommand{\selfcheckanswer}{\par\smallskip\textcolor{hsegray}{\bfseries Ответ.}\ }

\setlist[itemize]{itemsep=3pt,topsep=4pt}
\setlist[enumerate]{itemsep=3pt,topsep=4pt}
\setlength{\parindent}{0pt}
\setlength{\parskip}{5pt plus 1pt}
\setlength{\emergencystretch}{2em}
\setcounter{tocdepth}{2}

\pagestyle{fancy}
\fancyhf{}
\fancyhead[L]{\small\color{hsegray}Методы оптимизации}
\fancyhead[R]{\small\color{hsegray}Семинар 9}
\fancyfoot[L]{\small\color{hsegray}Евгений Баулин}
\fancyfoot[C]{\small\thepage}
\fancyfoot[R]{\small\color{hsegray}Лицензия MIT}
\renewcommand{\headrulewidth}{0pt}
\renewcommand{\footrulewidth}{0pt}
\fancypagestyle{plain}{\pagestyle{fancy}}

\begin{document}

\begin{center}
  {\large\color{hsegray}Методы оптимизации\par}
  \vspace{4pt}
  {\large\bfseries\color{hseaccent}Семинар 9\par}
  \vspace{8pt}
  {\LARGE\bfseries\color{hseblue}Линейное программирование: вершины, симплекс-метод, двойственность\par}
  \vspace{8pt}
  {\Large Конспект теории\par}
  \vspace{12pt}
  НИУ ВШЭ, факультет компьютерных наук\par
  2026-2027 учебный год\par
  \vspace{8pt}
  Лектор курса: Андрей Игнатов\par
  Евгений Баулин\par
  Telegram: \href{https://t.me/tarakan_tuc}{@tarakan\_tuc}\par
\end{center}

\noindent{\color{hseblue}\rule{\linewidth}{0.6pt}}

\tableofcontents

\clearpage


Этот конспект сопровождает семинар 9. Он начинается с геометрии линейной задачи, превращает вершины в базисные решения, строит на этом симплекс-метод и затем читает двойственное решение как список цен. Последние разделы показывают, как метод может сбоить и что всё это даёт целочисленной оптимизации. Все числа взяты из проверочного ноутбука семинара.

\begin{conventionbox}
Прямая задача: максимизировать $c^{\top}x$ при $Ax\le b$, $x\ge0$, где $A\in\mathbb R^{m\times n}$; $s=b-Ax$ — запасы по строкам (дополнительные). Двойственная: минимизировать $b^{\top}y$ при $A^{\top}y\ge c$, $y\ge0$. Последняя строка таблицы — $z-c^{\top}x=0$, поэтому отрицательный элемент отмечает переменную, которую стоит увеличивать. $B$ — базис, $x_B$ — базисные переменные, $p^*(b)$ — оптимум как функция правой части.
\end{conventionbox}

% =====================================================================
\section{Линейные задачи и их вершины}\label{sec:geometry}

Линейная задача оптимизирует линейную функцию на множестве, заданном конечным числом линейных неравенств и уравнений. Такое множество — многогранник: выпуклый, с плоскими гранями, возможно неограниченный, возможно пустой. Исходов у линейной задачи ровно три: она несовместна, неограниченна или имеет оптимальное решение, единственное или заполняющее целую грань, на плоскости — ребро.

\begin{definitionbox}{Определение: Стандартная форма.}
Задача $\min c^{\top}x$ при $Ax=b$, $x\ge0$. К ней сводится любая линейная задача: к строке $\le$ прибавляем дополнительную переменную, из строки $\ge$ вычитаем избыточную, свободную переменную записываем как $x^+-x^-$ с неотрицательными частями, строку с $b_i<0$ умножаем на $-1$, а максимизацию $c^{\top}x$ заменяем минимизацией $-c^{\top}x$.
\end{definitionbox}

\begin{definitionbox}{Определение: Базисное решение.}
Пусть у $A$ есть $m$ независимых строк. Выберем $m$ независимых столбцов, базис $B$, все прочие переменные приравняем к нулю и найдём $x_B$ из $Bx_B=b$. Получится базисное решение. Оно называется допустимым, если $x_B\ge0$, и вырожденным, если какая-то компонента $x_B$ равна нулю.
\end{definitionbox}

\begin{theorembox}{Теорема: Вершины — это базисные допустимые решения.}
Точка множества $\{x\ge0:Ax=b\}$ является вершиной, то есть не лежит посередине между двумя другими допустимыми точками, в точности тогда, когда она — базисное допустимое решение. Если у задачи в стандартной форме есть оптимум, то есть и оптимальная вершина.
\end{theorembox}

\begin{proofbox}
Если $x$ базисное и $x=\tfrac{1}{2}(u+v)$ с допустимыми $u,v$, то $u$ и $v$ обращаются в ноль там же, где $x$, и из $Bu_B=Bv_B=b$ следует $u=v=x$. Обратно, если столбцы при положительных компонентах $x$ зависимы, найдётся ненулевое $d$ на этом носителе с $Ad=0$, и точки $x\pm\varepsilon d$ допустимы, значит, $x$ не вершина. Об оптимуме: если у оптимальной $x$ носитель зависим, то $c^{\top}d=0$, иначе одна из точек $x\pm\varepsilon d$ была бы лучше $x$. Выберем знак $d$ так, чтобы у него была отрицательная компонента, и двигаемся вдоль него, пока какая-нибудь компонента точки не обнулится. Значение не меняется, носитель сокращается и через конечное число шагов становится независимым.
\end{proofbox}

\begin{whybox}
Линии уровня линейной функции плоские. Ничто не тянет оптимум внутрь множества, поэтому последнее касание всегда содержит вершину, и поиск можно ограничить конечным числом точек.
\end{whybox}

\handoutfigure{0.62}{fig_a1}{Смесь из A1. Закрашенные точки — допустимые базисные решения, крестики — базисные решения вне допустимого множества.}

\problemhead{A1}{Кормовая смесь}{У доски}
\textbf{Условие.} Смесь готовят из $x$ кг первого компонента по цене $5$ за кг и $y$ кг второго по цене $4$ за кг. В партии должно быть не меньше $8$ единиц белка и $9$ единиц клетчатки; в килограмме первого компонента белка $2$, а клетчатки $1$, в килограмме второго — $1$ и $3$ в тех же единицах. Минимизируйте стоимость графически, затем приведите задачу к стандартной форме и выпишите все базисные решения. Укажите $x^*$, стоимость $z^*$ и число допустимых базисных решений.

\textbf{Решение.}
\begingroup\sloppy
\begin{enumerate}[label=\textbf{\arabic*.}, leftmargin=*]
  \item Модель: $\min 5x+4y$ при $2x+y\ge8$, $x+3y\ge9$, $x,y\ge0$. Прямые пересекаются в точке $(3;\ 2)$, а допустимое множество — неограниченная область над обеими прямыми с угловыми точками $(0;\ 8)$, $(3;\ 2)$ и $(9;\ 0)$.
  \item Все линии уровня $5x+4y=z$ параллельны. Сдвигаем линию вниз к началу координат, пока она не готова потерять допустимое множество: последняя точка касания — угловая. Стоимости в угловых точках равны $(32;\ 23;\ 45)$, поэтому лучшая точка — $(3;\ 2)$ с $z^*=23$.
  \item В стандартной форме на каждую строку приходится одна избыточная переменная: $2x+y-s_1=8$, $x+3y-s_2=9$, все четыре переменные неотрицательны. Два уравнения и четыре неизвестных: выбираем две базисные переменные, остальные две обнуляем и решаем систему. Вариантов $\binom{4}{2}=6$.
  \item Все шесть разом:
\handoutdisplay{
\begin{array}{c|cccc|c|c} B & x & y & s_1 & s_2 & \ge 0 & z \\ \hline x,y & 3 & 2 & 0 & 0 & \checkmark & 23 \\ x,s_1 & 9 & 0 & 10 & 0 & \checkmark & 45 \\ x,s_2 & 4 & 0 & 0 & -5 & \times &  \\ y,s_1 & 0 & 3 & -5 & 0 & \times &  \\ y,s_2 & 0 & 8 & 0 & 15 & \checkmark & 32 \\ s_1,s_2 & 0 & 0 & -8 & -9 & \times &  \end{array}
}
 У трёх решений из шести все переменные неотрицательны, и это в точности три угловые точки. Остальные три — пересечения граничных прямых вне допустимого множества; отрицательная $s_i$ показывает, какое требование нарушено.
\end{enumerate}
\endgroup
\begin{answerbox}
$(x^*,y^*)=(3;\ 2)$, $z^*=23$, допустимых базисных решений 3 из 6.
\end{answerbox}

\begin{trapbox}
Слово «базисный» относится к выбору столбцов, а не к точке. Вырожденная вершина принадлежит нескольким базисам, а базисное решение может быть недопустимым. Считая базисы, вы не считаете вершины.
\end{trapbox}

\problemhead{A2}{Все базисные решения}{Самостоятельно}
\textbf{Условие.} Максимизируйте $3x+2y$ при $x+y\le7$, $x+2y\le10$, $x\le5$ и $x,y\ge0$. Добавьте три дополнительные переменные и переберите все наборы из трёх столбцов. Укажите, сколько наборов дают базисное решение, сколько из этих решений допустимы и чему равен оптимум.

\textbf{Решение.}
\begingroup\sloppy
\begin{enumerate}[label=\textbf{\arabic*.}, leftmargin=*]
  \item С дополнительными переменными система такова: $x+y+s_1=7$, $x+2y+s_2=10$, $x+s_3=5$. Три столбца из пяти выбираются $10$ способами. При наборе $\{y,s_1,s_2\}$ третья строка остаётся без базисной переменной, столбцы зависимы. Поэтому базисных решений $9$.
  \item Решаем каждую систему: неотрицательных решений $5$. Это вершины многоугольника $(0;\ 0)$, $(5;\ 0)$, $(5;\ 2)$, $(4;\ 3)$ и $(0;\ 5)$; у прочих какая-то дополнительная переменная отрицательна.
  \item Значения цели в этих вершинах: $(0;\ 15;\ 19;\ 18;\ 10)$. Наибольшее из них, $z^*=19$, достигается в $(5;\ 2)$.
\end{enumerate}
\endgroup
\begin{answerbox}
Базисных решений 9, допустимых среди них 5, $z^*=19$ в точке $(5;\ 2)$.
\end{answerbox}

\problemhead{A3}{Свободная переменная}{Самостоятельно}
\textbf{Условие.} Минимизируйте $x+2y$ при $x+y\ge1$, $x\le6$ и $x\ge0$, причём знак $y$ не ограничен. Запишите задачу в стандартной форме, заменив $y$ разностью двух неотрицательных переменных. Укажите число переменных в стандартной форме, затем $y^*$ и $z^*$.

\textbf{Решение.}
\begingroup\sloppy
\begin{enumerate}[label=\textbf{\arabic*.}, leftmargin=*]
  \item При $y=y^+-y^-$ система принимает вид $x+y^+-y^--s_1=1$, $x+s_2=6$: переменных $5$, это $x,y^+,y^-,s_1,s_2$, все они неотрицательны, а строк две. Цель: $x+2y^+-2y^-$.
  \item Цель уменьшается вместе с $y$, поэтому первая строка жёсткая: $y=1-x$, и $z=x+2(1-x)=2-x$. Эта величина убывает с ростом $x$, так что $x$ идёт к верхней границе.
  \item Итак, $x^*=6$, $y^*=-5$, $z^*=-4$. В стандартной форме $y^+=0$ и $y^-=5$: столбцы $y^+$ и $y^-$ противоположны, поэтому в один базис они не попадают.
\end{enumerate}
\endgroup
\begin{answerbox}
Переменных 5, $y^*=-5$, $z^*=-4$.
\end{answerbox}

\begin{selfcheckbox}
Почему столбцы $y^+$ и $y^-$ из A3 не могут быть базисными одновременно?
\selfcheckanswer Это противоположные векторы: они зависимы и вместе в базис входить не могут.
\end{selfcheckbox}

% =====================================================================
\section{Симплекс-метод}\label{sec:simplex}

Перебор базисов безнадёжен: число $\binom{n}{m}$ взрывается. Симплекс-метод смотрит на очень немногие из них. Он начинает в вершине и переходит к соседней с лучшим значением, пока лучших соседей не останется. Для линейной цели на выпуклом множестве вершина без лучших соседей — глобальный оптимум.

\begin{definitionbox}{Определение: Симплекс-таблица.}
Система $Ax+s=b$ вместе с $z-c^{\top}x=0$, переписанная так, что каждая базисная переменная стоит в своей строке одна. У каждой базисной переменной единичный столбец, последний столбец содержит их значения, а последняя строка показывает зависимость $z$ от небазисных переменных. Элементы $r_j$ этой строки — приведённые стоимости: $z=z_0-\sum_jr_jx_j$ по небазисным $j$.
\end{definitionbox}

\begin{itemize}
  \item Ввод: выбираем небазисный столбец с $r_j<0$, по правилу Данцига — наиболее отрицательный. Если такого нет, останавливаемся: точка оптимальна.
  \item Тест отношений: среди строк с положительным элементом в этом столбце находим наименьшее отношение правой части к элементу. Если положительных элементов нет, останавливаемся: задача неограниченна.
  \item Шаг: делим эту строку на разрешающий элемент и вычитаем её кратные из остальных строк, включая $z$-строку, чтобы столбец стал единичным.
\end{itemize}

\begin{theorembox}{Утверждение: Критерий оптимальности.}
Если в допустимом базисном решении все элементы $r_j$ строки $z$ неотрицательны, это решение оптимально.
\end{theorembox}

\begin{proofbox}
Строка $z=z_0-\sum_jr_jx_j$ — тождество на всех решениях системы. У любой допустимой точки $x_j\ge0$, поэтому её значение не больше $z_0$, а текущая точка достигает $z_0$. Если у небазисной переменной $r_j=0$, её ввод не меняет $z$, и если шаг ненулевой, оптимально всё ребро между двумя вершинами.
\end{proofbox}

\begin{whybox}
Тест отношений сохраняет допустимость: при росте вводимой переменной до $t$ каждая базисная переменная с положительным элементом $a$ уменьшается на $at$, и первая обнулившаяся останавливает движение. Она выходит из базиса, и мы приходим в соседнюю вершину.
\end{whybox}

\handoutfigure{0.62}{fig_b1}{Путь задачи B1: два шага от начала координат, вдоль оси $y$, а затем по верхнему ребру.}

\problemhead{B1}{Пекарня в таблице}{У доски}
\textbf{Условие.} Пекарня выпекает $x$ противней хлеба и $y$ противней пирогов; противень приносит $2$ и $5$. Время печи даёт $x+y\le8$, тесто — $2x+y\le14$, а начинка, нужная только пирогам, — $2y\le10$. Решите задачу симплекс-методом по правилу Данцига: в базис входит переменная с наиболее отрицательным элементом $z$-строки. Укажите $x^*$, $y^*$ и $z^*$.

\textbf{Решение.}
\begingroup\sloppy
\begin{enumerate}[label=\textbf{\arabic*.}, leftmargin=*]
  \item Добавляем $s_1,s_2,s_3$ и записываем $z-2x-5y=0$ последней строкой. Базис — дополнительные переменные, точка — начало координат, $z=0$. Рамкой отмечен первый разрешающий элемент.
\handoutdisplay{
\begin{array}{c|ccccc|c}  & x & y & s_1 & s_2 & s_3 & \\ \hline s_1 & 1 & 1 & 1 & 0 & 0 & 8 \\ s_2 & 2 & 1 & 0 & 1 & 0 & 14 \\ s_3 & 0 & \boxed{2} & 0 & 0 & 1 & 10 \\ \hline z & -2 & -5 & 0 & 0 & 0 & 0 \end{array}
}

  \item Наиболее отрицательный элемент стоит под $y$, значит, $y$ входит в базис: противень пирогов выгоднее всего. Насколько можно его увеличить? Отношения правых частей к положительным элементам столбца равны $(8;\ 14;\ 5)$, наименьшее — в третьей строке, поэтому $s_3$ выходит из базиса. Делим эту строку на разрешающий элемент и обнуляем остальной столбец.
\handoutdisplay{
\begin{array}{c|ccccc|c}  & x & y & s_1 & s_2 & s_3 & \\ \hline s_1 & \boxed{1} & 0 & 1 & 0 & -0{,}5 & 3 \\ s_2 & 2 & 0 & 0 & 1 & -0{,}5 & 9 \\ y & 0 & 1 & 0 & 0 & 0{,}5 & 5 \\ \hline z & -2 & 0 & 0 & 0 & 2{,}5 & 25 \end{array}
}

  \item Мы в точке $(0;\ 5)$, $z=25$. Отрицательный элемент остался только у $x$. Отношения равны $(3;\ 4{,}5)$, поэтому выходит $s_1$.
\handoutdisplay{
\begin{array}{c|ccccc|c}  & x & y & s_1 & s_2 & s_3 & \\ \hline x & 1 & 0 & 1 & 0 & -0{,}5 & 3 \\ s_2 & 0 & 0 & -2 & 1 & 0{,}5 & 3 \\ y & 0 & 1 & 0 & 0 & 0{,}5 & 5 \\ \hline z & 0 & 0 & 2 & 0 & 1{,}5 & 31 \end{array}
}

  \item В $z$-строке нет отрицательных элементов: $z=31-2s_1-\tfrac{3}{2}s_3$, и увеличение небазисной переменной может только уменьшить $z$. Значит, $(3;\ 5)$ — оптимум, $z^*=31$, а базисная $s_2=3$ показывает, сколько теста останется. Хватило двух шагов; если бы первым вошёл $x$, путь шёл бы в обход через $(7;\ 0)$ и $(6;\ 2)$.
\end{enumerate}
\endgroup
\begin{answerbox}
$(x^*,y^*)=(3;\ 5)$, $z^*=31$.
\end{answerbox}

\problemhead{B2}{Три переменные}{Самостоятельно}
\textbf{Условие.} Максимизируйте $5x_1+4x_2+3x_3$ при $x_1+x_3\le14$, $x_1+3x_2\le6$, $2x_1+x_2+x_3\le26$ и $x\ge0$. Решите задачу симплекс-методом с тем же правилом ввода. Укажите $z^*$, $x_3^*$ и итоговое значение третьей дополнительной переменной $s_3$.

\textbf{Решение.}
\begingroup\sloppy
\begin{enumerate}[label=\textbf{\arabic*.}, leftmargin=*]
  \item Начальная таблица, разрешающий элемент в рамке. Наибольший коэффициент у $x_1$, отношения $(14;\ 6;\ 13)$ выводят из базиса $s_2$.
\handoutdisplay{
\begin{array}{c|cccccc|c}  & x_1 & x_2 & x_3 & s_1 & s_2 & s_3 & \\ \hline s_1 & 1 & 0 & 1 & 1 & 0 & 0 & 14 \\ s_2 & \boxed{1} & 3 & 0 & 0 & 1 & 0 & 6 \\ s_3 & 2 & 1 & 1 & 0 & 0 & 1 & 26 \\ \hline z & -5 & -4 & -3 & 0 & 0 & 0 & 0 \end{array}
}

  \item После шага $z=30$. Теперь выгоден только $x_3$, отношения $(8;\ 14)$ выводят $s_1$.
\handoutdisplay{
\begin{array}{c|cccccc|c}  & x_1 & x_2 & x_3 & s_1 & s_2 & s_3 & \\ \hline s_1 & 0 & -3 & \boxed{1} & 1 & -1 & 0 & 8 \\ x_1 & 1 & 3 & 0 & 0 & 1 & 0 & 6 \\ s_3 & 0 & -5 & 1 & 0 & -2 & 1 & 14 \\ \hline z & 0 & 11 & -3 & 0 & 5 & 0 & 30 \end{array}
}

  \item Отрицательных элементов в $z$-строке нет, останавливаемся: $x^*=(6;\ 0;\ 8)$, $s_3=6$, $z^*=54$.
\handoutdisplay{
\begin{array}{c|cccccc|c}  & x_1 & x_2 & x_3 & s_1 & s_2 & s_3 & \\ \hline x_3 & 0 & -3 & 1 & 1 & -1 & 0 & 8 \\ x_1 & 1 & 3 & 0 & 0 & 1 & 0 & 6 \\ s_3 & 0 & -2 & 0 & -1 & -1 & 1 & 6 \\ \hline z & 0 & 2 & 0 & 3 & 2 & 0 & 54 \end{array}
}

\end{enumerate}
\endgroup
\begin{answerbox}
$x^*=(6;\ 0;\ 8)$, $z^*=54$, $s_3=6$.
\end{answerbox}

Методу нужна первая вершина. Для $Ax\le b$ при $b\ge0$ подходит начало координат. В остальных случаях делаем правые части неотрицательными, меняя знак строк, добавляем искусственную переменную в каждую строку без базисной переменной и на первой фазе минимизируем их сумму $w$ теми же шагами. Если минимум положителен, задача несовместна. Если он равен нулю, точка — вершина исходной задачи: искусственную переменную, оставшуюся в базисе на нуле, выводим шагом или удаляем её строку как избыточную, и отсюда начинается вторая фаза с настоящей целью.

\problemhead{B3}{Начало, которое надо найти}{Самостоятельно}
\textbf{Условие.} Максимизируйте $2x+y$ при $x+y\le6$, $x-y\ge2$ и $x,y\ge0$. Начало координат недопустимо, поэтому проведите первую фазу с одной искусственной переменной $a$ во второй строке, а затем вторую фазу с тем же правилом ввода. Укажите значение $x$ в конце первой фазы, оптимум $z^*$ и итоговое значение избыточной переменной $s_2$.

\textbf{Решение.}
\begingroup\sloppy
\begin{enumerate}[label=\textbf{\arabic*.}, leftmargin=*]
  \item Для строк $x+y+s_1=6$ и $x-y-s_2+a=2$ базис — $s_1,a$. Минимизировать $w=a$ значит максимизировать $-w$; исключив из этой строки базисную $a$, получаем
\handoutdisplay{
\begin{array}{c|ccccc|c}  & x & y & s_1 & s_2 & a & \\ \hline s_1 & 1 & 1 & 1 & 0 & 0 & 6 \\ a & \boxed{1} & -1 & 0 & -1 & 1 & 2 \\ \hline -w & -1 & 1 & 0 & 1 & 0 & -2 \end{array}
}

  \item Один шаг в столбце $x$ выводит $a$ и даёт $w=0$, значит, задача совместна, и первая фаза заканчивается в точке $(2;\ 0)$.
\handoutdisplay{
\begin{array}{c|ccccc|c}  & x & y & s_1 & s_2 & a & \\ \hline s_1 & 0 & 2 & 1 & 1 & -1 & 4 \\ x & 1 & -1 & 0 & -1 & 1 & 2 \\ \hline -w & 0 & 0 & 0 & 0 & 1 & 0 \end{array}
}

  \item Удаляем столбец $a$, возвращаем настоящую цель и исключаем базисную $x$: $z=4+3y+2s_2$.
\handoutdisplay{
\begin{array}{c|cccc|c}  & x & y & s_1 & s_2 & \\ \hline s_1 & 0 & \boxed{2} & 1 & 1 & 4 \\ x & 1 & -1 & 0 & -1 & 2 \\ \hline z & 0 & -3 & 0 & -2 & 4 \end{array}
}
 Затем входит $y$ и выходит $s_1$: единственный положительный элемент столбца стоит в её строке.
  \item Теперь $z=10$ в точке $(4;\ 2)$, а у $s_2$ всё ещё отрицательный элемент. В тесте отношений остаётся только первая строка, так что $y$ снова выходит.
\handoutdisplay{
\begin{array}{c|cccc|c}  & x & y & s_1 & s_2 & \\ \hline y & 0 & 1 & 0{,}5 & \boxed{0{,}5} & 2 \\ x & 1 & 0 & 0{,}5 & -0{,}5 & 4 \\ \hline z & 0 & 0 & 1{,}5 & -0{,}5 & 10 \end{array}
}

  \item Готово:
\handoutdisplay{
\begin{array}{c|cccc|c}  & x & y & s_1 & s_2 & \\ \hline s_2 & 0 & 2 & 1 & 1 & 4 \\ x & 1 & 1 & 1 & 0 & 6 \\ \hline z & 0 & 1 & 2 & 0 & 12 \end{array}
}
 Оптимум — $(6;\ 0)$, $z^*=12$, $s_2=4$. Вошедшая переменная может потом выйти, метод этого не запрещает.
\end{enumerate}
\endgroup
\begin{answerbox}
Первая фаза заканчивается в $(2;\ 0)$; $x^*=(6;\ 0)$, $z^*=12$, $s_2=4$.
\end{answerbox}

\begin{trapbox}
В тесте отношений отрицательные и нулевые элементы пропускаются. Деление на отрицательный элемент даст меньшее число, и если его взять, точка станет недопустимой.
\end{trapbox}

\begin{selfcheckbox}
В B1 в $z$-строке последней таблицы нули стоят под $x$, $y$ и $s_2$. Значит ли это, что оптимум не единственный?
\selfcheckanswer Нет. Это базисные переменные, а у них там всегда нули. Только ноль под небазисной переменной может указывать на другой оптимум.
\end{selfcheckbox}

% =====================================================================
\section{Двойственность}\label{sec:duality}

Откуда берётся верхняя оценка оптимума? Возьмём неотрицательные веса $y$ и сложим строки $Ax\le b$. Если взвешенная строка мажорирует цель, $A^{\top}y\ge c$, то $c^{\top}x\le y^{\top}Ax\le b^{\top}y$ для всех допустимых $x$. Поиск лучшей такой оценки сам является линейной задачей.

\begin{definitionbox}{Определение: Двойственная задача.}
Двойственная к $\max c^{\top}x$, $Ax\le b$, $x\ge0$ — это $\min b^{\top}y$, $A^{\top}y\ge c$, $y\ge0$. Строка прямой задачи даёт переменную двойственной, столбец — строку. Строка-равенство даёт свободную двойственную переменную, а свободная прямая переменная — строку-равенство. Двойственная к двойственной — снова прямая. Это двойственная по Лагранжу задача семинара 8, выписанная для линейной цели.
\end{definitionbox}

\begin{theorembox}{Теорема: Слабая и сильная двойственность.}
Для допустимых $x$ и $y$ выполнено $c^{\top}x\le b^{\top}y$. Если у одной из двух задач есть оптимальное решение, оно есть и у другой, и оптимальные значения равны. Если одна неограниченна, другая несовместна.
\end{theorembox}

\begin{proofbox}
Слабая двойственность — это цепочка выше. Для сильной доведём симплекс-метод до оптимальной таблицы с базисом $B$ и положим $y^{\top}=c_B^{\top}B^{-1}$. Элементы $z$-строки под дополнительными переменными равны $y$, а под $x_j$ — $(A^{\top}y-c)_j$; оптимальность говорит, что все они неотрицательны, значит, $y$ допустимо. Кроме того, $b^{\top}y=c_B^{\top}B^{-1}b=c^{\top}x$.
\end{proofbox}

\begin{theorembox}{Следствие: Дополняющая нежёсткость.}
Допустимые $x$ и $y$ оптимальны одновременно в точности тогда, когда $y_i(b-Ax)_i=0$ для каждой строки $i$ и $x_j(A^{\top}y-c)_j=0$ для каждого столбца $j$.
\end{theorembox}

\begin{proofbox}
Зазор равен $b^{\top}y-c^{\top}x=y^{\top}(b-Ax)+x^{\top}(A^{\top}y-c)$, это сумма неотрицательных произведений. Он равен нулю в точности тогда, когда равны нулю все произведения.
\end{proofbox}

\begin{whybox}
Это условия ККТ линейной задачи: допустимость прямой и двойственной задач и дополняющая нежёсткость, а стационарность заложена в строки двойственной задачи. Условие регулярности для линейной задачи не нужно.
\end{whybox}

\problemhead{C1}{Сколько стоит час работы печи}{У доски}
\textbf{Условие.} Возьмите пекарню из B1 с оптимумом $(3;\ 5)$. Запишите двойственную задачу, найдите $y^*$ из дополняющей нежёсткости и сравните с последней таблицей. При каких значениях предела времени печи $b_1$ базис остаётся оптимальным? Укажите $y_1^*$, $y_3^*$ и наибольшее такое $b_1$.

\textbf{Решение.}
\begingroup\sloppy
\begin{enumerate}[label=\textbf{\arabic*.}, leftmargin=*]
  \item Одна цена на ресурс, одна строка на продукт: минимизируем $8y_1+14y_2+10y_3$ при $y_1+2y_2\ge2$, $y_1+y_2+2y_3\ge5$, $y\ge0$. Это покупатель, арендующий все ресурсы: за каждый продукт он платит не меньше, чем тот принёс бы.
  \item Дополняющая нежёсткость в точке $(3;\ 5)$. Строка теста нежёсткая, $s_2=3$, поэтому $y_2=0$. И $x$, и $y$ положительны, значит, обе двойственные строки жёсткие: $y_1=2$, а из $y_1+2y_3=5$ следует $y_3=1{,}5$. Итак, $y^*=(2;\ 0;\ 1{,}5)$ — это элементы $z$-строки последней таблицы под $s_1,s_2,s_3$.
  \item Двойственное значение $b^{\top}y^*=31=z^*$, так что оптимальность обоих решений доказана сразу. Лишний час печи стоит $2$, лишнее тесто — ничего, единица начинки — $1{,}5$.
  \item Пусть $b_1=8+\delta$. Столбец $s_1$ последней таблицы показывает, как меняются базисные переменные: $x=3+\delta$, $s_2=3-2\delta$, $y=5$. Они неотрицательны при $-3\le\delta\le1{,}5$, то есть при $b_1$ из $[5,9{,}5]$. На этом отрезке наклон $p^*$ равен $y_1^*$; за его концами меняется базис, а с ним и цена.
\end{enumerate}
\endgroup
\begin{answerbox}
$y^*=(2;\ 0;\ 1{,}5)$, базис оптимален при $b_1\in[5,9{,}5]$.
\end{answerbox}

\problemhead{C2}{Решить вместо этого двойственную}{Самостоятельно}
\textbf{Условие.} Максимизируйте $9x_1+10x_2+3x_3$ при $2x_1+2x_2+5x_3\le14$, $3x_1+4x_2+2x_3\le24$ и $x\ge0$. У двойственной задачи две переменные: решите её графически, а затем восстановите решение прямой задачи из дополняющей нежёсткости. Укажите $y_1^*$, общее оптимальное значение и $x_1^*$.

\textbf{Решение.}
\begingroup\sloppy
\begin{enumerate}[label=\textbf{\arabic*.}, leftmargin=*]
  \item Двойственная задача: минимизировать $14y_1+24y_2$ при $2y_1+3y_2\ge9$, $2y_1+4y_2\ge10$, $5y_1+2y_2\ge3$, $y\ge0$.
  \item Её вершины — $(0;\ 3)$, $(3;\ 1)$ и $(5;\ 0)$ со значениями $(72;\ 66;\ 70)$. Наименьшее, $66$, достигается при $y^*=(3;\ 1)$. Третья строка здесь не работает: она ни в одной из вершин не активна.
  \item Дополняющая нежёсткость: при $y^*$ третья двойственная строка нежёсткая, поэтому $x_3=0$; оба $y_i>0$, значит, обе прямые строки жёсткие. Из $2x_1+2x_2=14$ и $3x_1+4x_2=24$ получаем $x^*=(4;\ 3;\ 0)$, и прямое значение, как и должно, равно $66$.
\end{enumerate}
\endgroup
\begin{answerbox}
$y^*=(3;\ 1)$, $x^*=(4;\ 3;\ 0)$, оптимальное значение $66$.
\end{answerbox}

% =====================================================================
\section{Теневые цены и диапазоны}\label{sec:sensitivity}

Сохраним оптимальный базис и немного изменим $b$. Базисные значения $x_B=B^{-1}b$ меняются линейно, и пока они неотрицательны, базис не перестаёт быть оптимальным: $z$-строка от $b$ не зависит. Значит, на этом диапазоне оптимум равен $y^{\top}b$ с тем же $y$.

\begin{theorembox}{Теорема: Теневые цены.}
Для невырожденного оптимального базиса $\partial p^*/\partial b_i=y_i^*$. В целом $p^*(b)$ вогнута и кусочно-линейна; каждый кусок отвечает одному базису, а в изломе левый и правый наклоны — две разные оптимальные цены.
\end{theorembox}

\begin{itemize}
  \item Диапазон $b_i$: столбец $s_i$ в последней таблице — это $i$-й столбец $B^{-1}$, поэтому $x_B+\delta B^{-1}e_i\ge0$ даёт отрезок для $\delta$.
  \item Диапазон $c_j$: точка та же, меняется $z$-строка. Пересчитайте её с новым коэффициентом и потребуйте, чтобы все элементы остались неотрицательными.
  \item Внутри диапазона $b$ цены те же, а план меняется; внутри диапазона $c$ план тот же, а цены меняются.
\end{itemize}

\handoutfigure{0.72}{fig_c1}{Оптимум задачи B1 как функция $b_1$. Наклон — теневая цена времени печи, в каждом изломе он падает.}

\problemhead{C3}{Докуда план выдерживает}{Самостоятельно}
\textbf{Условие.} Максимизируйте $4x+3y$ при $x+y\le9$, $2x+y\le14$, $y\le6$ и $x,y\ge0$. Найдите оптимум и цены. Затем укажите $y_1^*$, наибольшее $b_1$ (правая часть первой строки), при котором тот же базис остаётся оптимальным, и наибольший коэффициент при $x$ в цели, при котором та же точка остаётся оптимальной.

\textbf{Решение.}
\begingroup\sloppy
\begin{enumerate}[label=\textbf{\arabic*.}, leftmargin=*]
  \item Первые две строки пересекаются в $x^*=(5;\ 4)$, значение равно $32$, запас третьей строки — $2$. Цены находим из $y_1+2y_2=4$, $y_1+y_2=3$: $y^*=(2;\ 1;\ 0)$, цена третьей строки нулевая.
  \item При $b_1=9+\delta$ две жёсткие строки дают $x=5-\delta$ и $y=4+2\delta$. Одна только неотрицательность разрешила бы $\delta\le5$, но третья строка требует $y\le6$, то есть $\delta\le1$. Значит, базис держится до $b_1=10$, а дальше в дело вступает строка $y\le6$.
  \item Заменим коэффициент при $x$ на $c_1$. Точка остаётся оптимальной, пока линия уровня лежит между двумя жёсткими строками: их наклоны требуют $1\le c_1/3\le2$, то есть $c_1\in[3,6]$. То же следует из $y_1=2\cdot3-c_1\ge0$ и $y_2=c_1-3\ge0$.
\end{enumerate}
\endgroup
\begin{answerbox}
$y_1^*=2$, базис держится до $b_1=10$, точка — до коэффициента $6$.
\end{answerbox}

\begin{trapbox}
Теневая цена — производная, а не рыночная цена. За пределами диапазона она ничего не говорит, а в вырожденной вершине цены могут быть неединственными, и число может зависеть от солвера.
\end{trapbox}

\begin{selfcheckbox}
В C1 цена времени печи равна $2$ на $[5,9{,}5]$. Каковы наклоны $p^*$ ниже и выше этого отрезка?
\selfcheckanswer Ниже — $5$, потому что тогда печь ограничивает одни пироги; выше — $0$, потому что печи уже хватает.
\end{selfcheckbox}

% =====================================================================
\section{Вырожденность, зацикливание и время работы}\label{sec:failures}

Конечность симплекс-метода опирается на один факт: каждый шаг строго увеличивает $z$, поэтому базисы не повторяются. Это перестаёт быть верным, когда базисная переменная равна нулю: тест отношений может дать шаг нулевой длины. Такой вырожденный шаг меняет базис, но не двигает точку.

\begin{theorembox}{Теорема: Правило Бленда.}
Если среди столбцов с отрицательным элементом вводить переменную с наименьшим индексом, а среди строк с наименьшим отношением выводить переменную с наименьшим индексом, симплекс-метод останавливается за конечное число шагов.
\end{theorembox}

На практике это правило медленное, поэтому программы используют его как запасное или слегка возмущают $b$, чтобы разбить ничьи. Что до скорости, куб Кли — Минти показывает, что правило Данцига может пройти все $2^n$ вершин, и ни одно правило выбора не известно как полиномиальное. Но на настоящих задачах шагов обычно лишь в небольшое число раз больше, чем $m$. Полиномиальные алгоритмы для ЛП существуют: это метод эллипсоидов и методы внутренней точки, причём последние на практике соперничают с симплекс-методом.

\handoutfigure{0.72}{fig_km}{Число шагов правила Данцига на кубах Кли — Минти и на случайных задачах тех же размеров.}

\begin{itemize}
  \item \texttt{linprog}, status 3: неограниченна. Ищите забытое ограничение.
  \item \texttt{linprog}, status 2: несовместна. Ищите противоречащие друг другу строки; сертификат из леммы Фаркаша их указывает.
  \item \texttt{res.ineqlin.marginals} содержит цены задачи минимизации: если задачу развернули сменой знака, у них тоже меняется знак.
\end{itemize}

% =====================================================================
\section{Зачем всё это целочисленной оптимизации}\label{sec:integer}

Солвер целочисленных моделей тратит большую часть времени на линейные задачи. Релаксация семинара 8 — это линейная задача, а её значение — оценка, с которой сравнивается каждая целая точка. Сильная двойственность делает эту оценку дешёвой для проверки: любое допустимое $y$ даёт её без доверия к солверу.

Таблица тоже помогает. Новая граница на переменную или новая верная строка оставляют старый базис двойственно допустимым, поэтому следующая задача стартует оттуда, обычно за несколько шагов, а не с нуля. А если оптимальная вершина релаксации оказалась целой, больше делать нечего: для этой подзадачи поиск окончен.

\end{document}
